\documentclass[10pt]{article}
\usepackage{icmlko}

\icmlmeta{IP 파운데이션 월드모델}{151}{필터인 줄 알았는데 사람 손이었다}{팝업 유보 59건이 모든 판정의 병목이라 그 문을 열러 갔다 --- 문 뒤에 있던 것은 집계 방식이 아니라 축을 매긴 사람이었다}

\begin{document}

\twocolumn[
\icmlheader

\begin{abstract}
노트 148이 ``남은 길은 축과 표본''이라 했고 축은 노트 149 · 150에서 위키로
$+$0.0104를 벌었다. 이번은 표본이다. \textbf{팝업 유보 59건이 이 판의 모든
판정을 막는다} --- 짝 반폭이 $1.03\,n^{-0.52}$로 0.124라, 노트 148이 잰
``판 순위와 팝업 순위가 무관하다''도 팝업 쪽에서는 확인할 수가 없다. 어디서
막히나를 갈랐다. \textbf{필터 하나였다} --- 원본 376건이 라벨(350) · 등급
A·B(193) · 범위(189)를 지나 \textbf{집계 방식 필터에서 75건}이 된다. 그것만
풀면 유보가 59에서 116으로 는다. 반폭이 0.124에서 0.087로 좋아진다.
\textbf{그런데 늘어난 57건은 신호가 없다.} 정식화 셋이 전부 음수다 ---
F18 배깅 $-$0.085 · F8 부스팅 $-$0.088 · F6 능형 $-$0.073. 116건 전체로
재면 팝업 $\rho$가 0.43에서 0.35로 내려간다. \textbf{분모만 늘려 분해능을
사는 척한 것이다}(노트 127의 묶음 함정과 같은 모양). 여기서 멈추지 않고
왜 그런지 갈랐다. 집계 방식별로 보니 \texttt{organizer\_claim} 45건의
$\rho$가 \emph{정의되지 않는다} --- 모형이 그 45건에 전부 같은 점수를 준다.
이유는 하나다. \textbf{그 45건에는 축이 안 매겨져 있다.} 집계 방식별 축
빈칸 비율을 재니 \texttt{organizer\_claim} 이 \textbf{98/98 = 100\%}이고
표준 집계는 10.7\%다. \textbf{집계 필터는 집계 품질을 거른 것이 아니라
사람이 축을 매겼는지를 거르고 있었다} --- 두 조건이 거의 완전히 얽혀 있다.
그러므로 팝업의 천장은 필터가 아니라 \textbf{사람이 태깅한 건수}다. 노트
132가 잰 상한 202도, 노트 133이 자동 태깅으로 팝업을 $+$0.4562에서
$+$0.1150으로 무너뜨린 것도 같은 벽의 다른 면이었다.
\end{abstract}
]

\begin{figure*}[t]
\centering
\includegraphics[width=\textwidth]{figs/gate.pdf}
\caption{왼쪽 --- 한 필터가 189를 75로 줄인다. 가운데 --- 늘어난 절반은
정식화 셋 다 음수다. 오른쪽 --- 그 필터는 사람 손의 대리였다.}
\end{figure*}

\section{어디서 막히나}

\begin{table}[t]
\centering\tsize
\caption{팝업 376건이 59건이 되는 길.}
\begin{tabular}{lrr}
\toprule
& 남음 & 잃음 \\
\midrule
원본 & 376 & \\
라벨 있음 & 350 & 26 \\
등급 A·B & 193 & 157 \\
범위 쓸 수 있음 & 189 & 4 \\
\textbf{집계 방식 표준} & \textbf{75} & \textbf{114} \\
\midrule
(그중 2025 이후) & \textbf{59} & \\
\bottomrule
\end{tabular}
\end{table}

\claimx{\textbf{한 필터가 절반 넘게 가져간다.} 집계 방식 필터
(\texttt{entry} · \texttt{participation} 만 통과)가 189를 75로 줄인다.
그것만 풀면 유보가 59에서 116으로 늘고 짝 반폭이 0.124에서 0.087로
내려간다.}{푸는 것 자체는 노트 133이 이미 학습 쪽에서 했다 --- 넓힌 팝업.
그때도 \emph{채점} 분모는 59로 뒀다. 왜 그랬는지가 이 노트의 답이다.}

\section{늘어난 절반은 신호가 없다}

\begin{table}[t]
\centering\tsize
\caption{같은 모형, 유보를 갈라 잰 팝업 $\rho$.}
\begin{tabular}{lrrr}
\toprule
& 전체 116 & 표준 59 & \textbf{늘어난 57} \\
\midrule
F6 능형 & $+$.458 & $+$.459 & \textbf{$-$.073} \\
F8 부스팅 & $+$.316 & $+$.372 & \textbf{$-$.088} \\
F18 배깅 & $+$.352 & $+$.432 & \textbf{$-$.085} \\
\bottomrule
\end{tabular}
\end{table}

\claimx{\textbf{정식화 셋이 전부 음수다.} 서로 아주 다른 모형(선형 · 부스팅 ·
배깅)이 같은 방향으로 0에 못 미친다. 모형 탓이 아니라 그 57건에 잴 것이
없다는 뜻이다.}{반폭은 정말로 좋아진다(0.124$\to$0.087). \textbf{분모를
늘려 분해능을 사면 안 된다} --- 노트 127에서 축 없는 51건을 더해 팝업이
$+$0.3662에서 $+$0.4459로 ``올랐던'' 것과 같은 자리다.}

\section{왜 그런가 --- 갈라 봤다}

\begin{table}[t]
\centering\tsize
\caption{집계 방식별. 유보 116건.}
\begin{tabular}{lrrl}
\toprule
& n & $\rho$ & \\
\midrule
참여(participation) & 19 & $+$.477 & 표준 \\
입장(entry) & 40 & $+$.346 & 표준 \\
\textbf{주최측 주장} & \textbf{45} & \textbf{정의 안 됨} & 비표준 \\
불명 & 10 & --- & 표본 부족 \\
\bottomrule
\end{tabular}
\end{table}

\claimx{\textbf{$\rho$ 가 정의되지 않는다 --- 라벨이 아니라 예측 때문이다.}
그 45건의 라벨은 34가지 서로 다른 값이다. 정의가 안 되는 쪽은 \emph{예측}
이다 --- 모형이 45건에 전부 같은 점수를 준다.}{같은 점수를 준다는 것은
모형이 그 45건을 구별할 재료가 없다는 뜻이다. 무엇이 없나.}

\begin{table}[t]
\centering\tsize
\caption{집계 방식별 ``축이 안 매겨진'' 비율.}
\begin{tabular}{lrr}
\toprule
& n & 축 빈칸 \\
\midrule
입장 & 46 & 8.7\% \\
참여 & 29 & 13.8\% \\
불명 & 14 & 14.3\% \\
\textbf{주최측 주장} & \textbf{98} & \textbf{100.0\%} \\
\midrule
표준 전체 & 75 & 10.7\% \\
\textbf{비표준 전체} & \textbf{114} & \textbf{87.7\%} \\
\bottomrule
\end{tabular}
\end{table}

\claimx{\textbf{집계 필터는 사람 손의 대리였다.} \texttt{organizer\_claim}
98건이 하나도 빠짐없이 축이 비어 있다. 표준 집계는 10.7\%만 비어 있다.
두 조건이 거의 완전히 얽혀 있어서, ``집계 방식으로 걸렀다''고 적어 온 것이
사실은 ``사람이 태깅한 것만 남겼다''였다.}{인과가 반대일 수도 있다 ---
주최측 주장만 있는 건은 자료가 얇아서 사람이 태깅을 포기했을 수 있다.
어느 쪽이든 \emph{같이 움직인다}는 사실은 바뀌지 않는다.}

\section{그래서 천장이 무엇인가}

\begin{table}[t]
\centering\tsize
\caption{팝업 유보를 늘리는 세 길.}
\begin{tabular}{lrl}
\toprule
& 유보 & \\
\midrule
현행 & 59 & \\
집계 필터 풀기 & 116 & \textbf{절반이 잡음 --- 안 됨} \\
그 45건을 사람이 태깅 & 116 & 과제 78의 규모 \\
자동 태깅 & 116 & 노트 133: $+$.456$\to$$+$.115 \\
\bottomrule
\end{tabular}
\end{table}

\claimx{\textbf{천장은 필터가 아니라 태깅 건수다.} 노트 132가 잰 상한 202,
노트 133이 자동 태깅으로 팝업을 무너뜨린 것, 그리고 이 노트의 100\% 빈칸 ---
셋이 같은 벽의 세 면이다.}{그러면 과제 78(시장 레코드 20건 사람 태깅으로
자동 태거의 전이 오차 측정)이 부수 과제가 아니라 \textbf{임계 경로}다.
자동 태거가 쓸 만해지면 유보가 두 배가 되고 반폭이 0.124에서 0.087로 준다.}

\section{덧 --- 상태 층을 찾아봤다}

같은 사이클에서 BQ 를 읽기 전용으로 훑었다. 지금 라벨은 스칼라 하나(일평균)
이고 월드모델이라면 궤적이어야 하므로, 일별 자료가 있는지 봤다.

\begin{table}[t]
\centering\tsize
\caption{BQ 일별 표. 읽기만 했다.}
\begin{tabular}{lrl}
\toprule
& 행 & \\
\midrule
core.popup\_visitor\_daily & 0 & 비어 있다 \\
mart.popup\_daily\_sales & 0 & 비어 있다 \\
core.tenant\_sales\_daily & 0 & 비어 있다 \\
\textbf{core.retail\_sale\_daily} & \textbf{6{,}636} & \textbf{브랜드 187} \\
\bottomrule
\end{tabular}
\end{table}

\claimx{\textbf{궤적 자료가 하나 있다.} \texttt{retail\_sale\_daily} 는
2025-03부터 2026-06까지 브랜드 187개의 일별 GMV · 거래 건수다. 7일 이상이
171개, 31일 이상이 42개다. \textbf{171 궤적은 팝업 유보 59 스칼라보다
크다.}}{다만 \emph{다른 모집단}이다. \texttt{engagement.primary\_brand\_id}
가 572건 전부 NULL 이고 브랜드 이름과 프로젝트 이름의 겹침이 0이다 ---
이쪽은 매장 입점 브랜드(카페 · 베이커리)이고 팝업 레코드는 대행 프로젝트다.
축을 새로 매겨야 쓸 수 있다.}

\begin{table}[t]
\centering\tsize
\caption{남은 것.}
\begin{tabular}{ll}
\toprule
\textbf{과제 78} & 임계 경로가 됐다 --- 자동 태거 전이 오차 \\
상태 층 & 브랜드 171 궤적 --- 축을 새로 매겨야 한다 \\
집계 & 얽힘의 인과 방향은 아직 모른다 \\
\bottomrule
\end{tabular}
\end{table}

첫 줄이 다음이다. 지금까지 과제 78을 ``해 두면 좋은 것''으로 미뤄 왔는데,
이 노트가 그것을 병목 위에 올려놨다. \textbf{자동 태거의 전이 오차를 재기
전에는 팝업 표본을 늘릴 방법이 없다.}

\end{document}
